\documentclass[12pt,a4paper]{article}

% Настройки шрифтов и локализации для XeLaTeX
\usepackage{fontspec}
\usepackage{polyglossia}
\setmainlanguage[babelshorthands=false]{russian}
\setotherlanguage{english}

% Попытка установить стандартные шрифты (можно заменить на любые системные, например, Times New Roman)
\setmainfont{FreeSerif}
\setsansfont{FreeSans}
\setmonofont{FreeMono}

% Математические пакеты
\usepackage{amsmath, amsthm, amssymb}
\usepackage{mathtools}
\usepackage{calrsfs}
\usepackage{enumitem}
\usepackage{microtype}
\usepackage{geometry}
\usepackage{hyperref}
\usepackage{xcolor}

% Паттерны переноса для русского языка недоступны в этом окружении offline
% (hyph-ru отсутствует в установке TeX Live, сеть отключена). Поэтому вместо
% переноса слов разрешаем абзацному алгоритму растягивать межсловные пробелы
% и слегка нарушать «жёсткость» строки — это предотвращает выход текста
% за поля страницы (Overfull \hbox) ценой чуть менее плотного набора.
\tolerance=2000
\emergencystretch=3em
\hbadness=10000
\hfuzz=2pt
\sloppy

\geometry{
  a4paper,
  top=2.5cm, bottom=2.5cm,
  left=3cm, right=3cm
}

\hypersetup{
  colorlinks=true,
  linkcolor=blue!60!black,
  urlcolor=blue!60!black
}

% Пакеты для графики и диаграмм
\usepackage{tikz}
\usetikzlibrary{cd, arrows.meta, shapes.geometric}

% Определение математических окружений
\theoremstyle{plain}
\newtheorem{theorem}{Теорема}
\newtheorem{lemma}[theorem]{Лемма}
\newtheorem{proposition}[theorem]{Предложение}
\newtheorem{corollary}[theorem]{Следствие}

\theoremstyle{definition}
\newtheorem{definition}[theorem]{Определение}
\newtheorem{remark}[theorem]{Замечание}
\newtheorem{example}[theorem]{Пример}

% Математические макросы для унификации обозначений пользователя
\newcommand{\ph}{\varphi}
\newcommand{\toinf}{\rightarrow}
\newcommand{\PA}{\textsf{PA}}
\newcommand{\Ztwo}{\textsf{Z}_2}
\newcommand{\ZFC}{\textsf{ZFC}}
\newcommand{\PArec}{\PA^{\mathrm{rec}}_{-}}
\newcommand{\eqdef}{\mathrel{\overset{\mathrm{def}}{\rightleftharpoons}}}

\begin{document}

\title{\textbf{О математической индукции}}
\author{Н.\,И.\,Казимиров}
\date{2026}
\maketitle
\tableofcontents
\bigskip

\section{Формализации и разновидности индукции}

Принцип математической индукции представляет собой несущую конструкцию регулярных дискретных систем. В зависимости от экспрессивной мощности языка и дедуктивной силы базовой теории, индукция принимает существенно различные логические формы.

\subsection{Первопорядковая и второпорядковая формулировки}
В первопорядковой арифметике Пеано ($\PA$) индукция не может быть выражена единой формулой в силу невозможности квантификации по предикатам: формально она постулируется как счётная схема аксиом, по одной для каждой формулы $\ph(a)$ (возможно, с параметрами $\vec p$):
$$\ph[a/0] \land \forall x (\ph[a/x] \toinf \ph[a/Sx]) \toinf \forall y \; \ph[a/y].$$
Содержательно, однако, эту схему привычнее и понятнее записывать в форме, к которой математики привыкли по индуктивным определениям в теории множеств, --- отождествляя формулу $\ph$ с её объёмом $X=\{x \mid \ph(x)\}$:
\begin{equation}
    I:\qquad \forall X \Big(0\in X \,\land\, \forall x\,(x\in X \toinf Sx\in X) \;\toinf\; \forall x\,(x\in X)\Big).
\end{equation}
Нужно сразу же оговорить, в каком смысле следует понимать запись $\forall X$ в (1): в первопорядковом языке $\PA$ квантификации по множествам не существует, так что (1) --- это \emph{метазапись} схемы аксиом (по всем формулам $\ph$, задающим $X$), а не одна формула с настоящим квантором второго порядка. Именно эту схему -- обозначаемую $I$ в согласии с уже принятыми обозначениями $I\Sigma_n$, $I\Delta_0$ для её фрагментов по сложности формулы $\ph$ -- мы подробно обсудим в §2, вместе с несколькими равносильными ей принципами.

В системе анализа второго порядка ($\Ztwo$), где квантификация по множествам разрешена самим языком, схема (1) сворачивается в буквальную единую аксиому того же вида -- уже без оговорки о метазаписи:
\begin{equation}
    \forall X \left( 0 \in X \land \forall x (x \in X \toinf S(x) \in X) \toinf \forall y (y \in X) \right).
\end{equation}
Дедуктивная мощность $\Ztwo$ определяется тем, какая \textit{схема свёртывания} (comprehension) постулируется для существования множеств $X = \{x \mid \ph(x)\}$: чем шире класс формул $\ph$, для которых гарантировано существование такого $X$, тем на больший запас «множеств» распространяется единая аксиома (2). При непредикативном свёртывании (принятом в полной $\Ztwo$) это охватывает уже произвольное подмножество, определимое во всём языке анализа.

\subsection{Разновидности индукции по сложности формул и параметрам}
\begin{enumerate}[label=\arabic*.]
    \item \textbf{Ограничение по иерархии формул ($I\Sigma_n$, $I\Pi_n$):} Ограничение кванторной сложности формулы $\ph$ в схеме индукции порождает строгую дедуктивную иерархию подсистем арифметики. Например, $I\Sigma_1$ достаточна для доказательства тотальности примитивно рекурсивных функций, тогда как $I\Delta_0$ (индукция по формулам с ограниченными кванторами) не способна доказать тотальность экспоненты.
    \item \textbf{Индукция с параметрами и без (Parameter-free induction):} В полной $\PA$ наличие свободных переменных $\vec{p}$ в формуле $\ph$ не увеличивает силу теории. Однако для ограниченных фрагментов типа $I\Sigma_n^-$ (без параметров) сила теории строго меньше, чем для фрагментов с параметрами $I\Sigma_n$.
    \item \textbf{Индукция вдоль упорядочений:} Обобщение локального шага на произвольные вполне упорядоченные множества.
\end{enumerate}

\subsection{Индукция в теории множеств}
В теории множеств Цермело--Френкеля ($\ZFC$) индукция по натуральным числам является следствием аксиомы бесконечности, гарантирующей существование индуктивного множества $\omega$:
\begin{equation}
    \exists X \left( \emptyset \in X \land \forall x \in X (x \cup \{x\} \in X) \right).
\end{equation}
Множество $\omega$ определяется как пересечение всех таких индуктивных множеств, из чего тривиально следует стандартная индукция для подмножеств $\omega$.

\textit{Трансфинитная индукция} обобщает этот принцип на класс всех ординалов $\text{On}$ (или на произвольные отношения фундированности):
\begin{equation}
    \forall \alpha \in \text{On} \left( (\forall \beta < \alpha \; \ph(\beta)) \toinf \ph(\alpha) \right) \toinf \forall \gamma \in \text{On} \; \ph(\gamma).
\end{equation}
В $\ZFC$ трансфинитная индукция как принцип доказательства (в отличие от трансфинитной \emph{рекурсии}, для определения функций которой существенна именно схема подстановки) является теоремой, опирающейся лишь на аксиому регулярности (фундированности) и схему выделения.

\section{Равносильные первопорядковые принципы и метаматематические тонкости}

\subsection{Пять принципов: определения}
Рассмотрим пять ключевых принципов, формулируемых в первопорядковом языке арифметики. Как и в §1.1, каждый из них -- схема аксиом (по формуле $\ph$, задающей $X=\{x \mid \ph(x)\}$), и запись $\forall X(\ldots)$ всюду означает именно эту схему, а не квантор второго порядка. Для $x$ обозначим через $[0,x)\eqdef\{y \mid y<x\}$ начальный интервал, и договоримся о трёх сокращениях:
$$\mathrm{Ind}(X) \eqdef\; 0\in X \,\land\, \forall x\,(x\in X\toinf Sx\in X),$$
$$\mathrm{Prog}(X) \eqdef\; \forall x\,\big([0,x)\subseteq X\;\toinf\; x\in X\big),$$
$$\mathrm{Tot}(X) \eqdef\; \forall x\,(x\in X).$$
Здесь \emph{индуктивность} -- условие локального шага; \emph{прогрессивность} -- замкнутость относительно начальных интервалов; \emph{тотальность} -- утверждение, что $X$ содержит вообще все элементы. Тогда:
\begin{enumerate}[leftmargin=1.2em, itemindent=0pt, labelsep=0.4em, itemsep=0.7em]
    \item \textbf{Локальная индукция ($I$).} \\ $\forall X\big(\mathrm{Ind}(X)\toinf \mathrm{Tot}(X)\big)$.
    \item \textbf{Возвратная (полная) индукция ($I^*$).} \\ $\forall X\big(\mathrm{Prog}(X)\toinf\mathrm{Tot}(X)\big)$.
    \item \textbf{Принцип наименьшего элемента ($L$).} \\ $\forall X\big(X\ne\varnothing \;\toinf\; \exists x\in X\,\forall y\in X\,(x\leqslant y)\big)$.
    \item \textbf{Принцип бесконечного спуска ($D$).} \\ $\forall X\big(X\ne\varnothing \;\toinf\; \neg\big(\forall x\in X\,\exists y\in X\,(y<x)\big)\big)$.
    \item \textbf{Схема коллекции ($B$).} \\ $\forall a\,\big(\forall x{<}a\,\exists y\,\ph(x,y) \;\toinf\; \exists z\,\forall x{<}a\,\exists y{<}z\,\ph(x,y)\big)$.
\end{enumerate}
(Схему $B$ неудобно и неестественно записывать через характеристическое множество $X$ -- по своей форме она устроена иначе, как утверждение о двуместном отношении $\ph(x,y)$, -- поэтому для неё мы сохраняем обычную формульную запись.)

\subsection{Равносильность $I^*$, $L$, $D$ в чистом исчислении предикатов}
Принципиально важное наблюдение: $I^*$, $L$ и $D$ равносильны друг другу уже в чистом исчислении предикатов -- без единой арифметической аксиомы, -- если потребовать от $<$ быть отношением строгого линейного порядка. Нам не нужна здесь подробная аксиоматика порядка (это не наша цель); важно лишь одно её следствие -- \emph{трихотомия}: для любых $x,y$ верно ровно одно из трёх: $x<y$, $x=y$, $y<x$. Трихотомия используется ровно один раз: чтобы связать «отрицательную» форму принципа $D$ (выраженную через $<$) с его логическим двойником -- привычной для математика «позитивной» формой $L$ (через $\leqslant$, сравнение с \emph{каждым} элементом $X$).

\begin{proposition}
Для любой формулы $\ph$ (с $X=\{x \mid \ph(x)\}$):
\begin{enumerate}[label=(\alph*)]
    \item в чистом исчислении предикатов $I^*_{\neg\ph} \iff D_\ph$;
    \item если, кроме того, $<$ удовлетворяет трихотомии, то $D_\ph \iff L_\ph$.
\end{enumerate}
Поскольку класс формул замкнут относительно $\ph\mapsto\neg\ph$, отсюда вытекает равносильность полных схем $I^*\iff D\iff L$.
\end{proposition}
\begin{proof}
(a) Схема $I^*$ для формулы $\neg\ph$ имеет вид $A\toinf B$, где (разворачивая $\mathrm{Prog}$ и $\mathrm{Tot}$ по определению)
$$A \equiv \forall x\big(\forall y{<}x\,(\neg\ph(y)) \toinf \neg\ph(x)\big), \qquad B\equiv \forall x\,(\neg\ph(x)).$$
По закону контрапозиции $A\toinf B$ равносильно $\neg B\toinf\neg A$ -- чистая пропозициональная тавтология, не требующая никаких новых аксиом. Раскрывая отрицания по правилам де Моргана и двойственности кванторов:
$$\neg B \equiv \exists x\,(\ph(x)) \equiv X\ne\varnothing, \qquad \neg A \equiv \exists x\big(\ph(x)\,\land\,\forall y{<}x\,(\neg\ph(y))\big).$$
Внутри $\neg A$ подформула «$\forall y{<}x\,\neg\ph(y)$», то есть «$\forall y\,(y<x\toinf\neg\ph(y))$», сама по чистой логике равносильна своей контрапозиции «$\forall y\,(\ph(y)\toinf\neg(y<x))$», то есть «$\forall y\in X\,\neg(y<x)$». Значит,
$$\neg A \equiv \exists x\in X\,\forall y\in X\,\neg(y<x) \equiv \neg\big(\forall x\in X\,\exists y\in X\,(y<x)\big),$$
и импликация $\neg B\toinf\neg A$ -- это в точности $D_\ph$.

(b) Раскрыв внутреннее отрицание, $D_\ph$ утверждает: $X\ne\varnothing\toinf\exists x\in X\,\forall y\in X\,\neg(y<x)$. Формула $L_\ph$ устроена буквально так же, только с «$\neg(y<x)$», заменённым на «$x\leqslant y$». По трихотомии это одно и то же условие: $\neg(y<x)$ означает, что либо $y=x$, либо $x<y$, то есть ровно $x\leqslant y$.
\end{proof}

Итак, равносильность $I^*\iff L\iff D$ -- факт \emph{логики}, а не арифметики: доказательство не использует ни число $0$, ни операцию $+1$, а лишь линейность порядка $<$.

\subsection{Равносильность $I$ и $I^*$: роль минимальной арифметики последователя}
Связь локальной индукции $I$ с тройкой $I^*, L, D$ устроена принципиально иначе: она неотделима от последователя $S$ и потому не может быть чисто логическим фактом. Нам понадобятся три элементарных факта о $<$ и $S$. Здесь стоит быть аккуратным с тем, откуда они берутся: в арифметике Робинсона $\mathsf{Q}$ отношение $<$ обычно вообще не примитивно, а \emph{определяется} экзистенциальной формулой ($x<y \eqdef \exists z\,(Sz+x=y)$), и доказательство того, что это определение равносильно приводимым ниже рекурсивным тождествам, само по себе не совсем тривиально и зависит от способа рекурсивного разворачивания $+$. Чтобы не оставлять это висящим в воздухе, зафиксируем $<$ как примитивное отношение и введём теорию
$$\PArec \;\eqdef\; \big\{\,x+0=x,\ x+Sy=S(x+y),\ x\cdot0=0,\ x\cdot Sy=x\cdot y+x,$$
$$Sx\ne0,\ Sx=Sy\to x=y\,\big\} \,\cup\, \{\text{(O1)--(O3)}\}$$
-- по существу это арифметика Робинсона в её исходной, рекурсивной форме, с явно постулированными (а не выведенными) аксиомами порядка (O1)--(O3):
\begin{itemize}
    \item[(O1)] $\neg(x<0)$ -- $0$ есть наименьший элемент;
    \item[(O2)] $y<Sx \leftrightarrow y\leqslant x$ -- рекурсивная связь порядка с последователем;
    \item[(O3)] $x\ne0 \toinf \exists y\,(x=Sy)$ -- всякий ненулевой элемент есть последователь.
\end{itemize}
Все три -- аксиомы $\PArec$, а не теоремы, выводимые из чего-то более простого.

\begin{lemma}
Из (O1)--(O3) следует: для всякого $X$, $\mathrm{Ind}(X)\toinf\mathrm{Prog}(X)$.
\end{lemma}
\begin{proof}
Пусть выполнено $\mathrm{Ind}(X)$, то есть $0\in X$ и $\forall z\,(z\in X\toinf Sz\in X)$. Возьмём произвольный $x$ с $[0,x)\subseteq X$; покажем $x\in X$. По (O3) либо $x=0$, либо $x=Sz$ для некоторого $z$.
Если $x=0$: тогда $x\in X$ прямо по первому конъюнкту $\mathrm{Ind}(X)$.
Если $x=Sz$: то по (O2) (при $y=z$) имеем $z<Sz=x$, откуда $z\in[0,x)\subseteq X$, то есть $z\in X$; тогда по второму конъюнкту $\mathrm{Ind}(X)$ получаем $Sz=x\in X$.
\end{proof}

\begin{proposition}
При наличии (O1)--(O3), из $I^*$ следует $I$.
\end{proposition}
\begin{proof}
Пусть выполнено $\mathrm{Ind}(X)$. По Лемме, $\mathrm{Prog}(X)$. Применяя $I^*$, получаем $\mathrm{Tot}(X)$ -- заключение $I$.
\end{proof}
Отметим: аксиома (O3) использована здесь единственный раз -- в разборе случаев для $x$ (либо $0$, либо последователь). Без неё нельзя было бы гарантировать, что каждый ненулевой элемент <<достижим шагом индукции>> из какого-то предыдущего.

\begin{proposition}
При наличии лишь (O1)--(O2) (аксиома (O3) не требуется!), из $I$ следует $I^*$.
\end{proposition}
\begin{proof}
Пусть выполнено $\mathrm{Prog}(X)$. Определим вспомогательное множество
$$Y \eqdef \{x \mid [0,x)\subseteq X\}$$
-- максимальный начальный интервал, целиком лежащий в $X$.

\emph{$Y\subseteq X$.} Если $x\in Y$, то $[0,x)\subseteq X$ -- а это в точности посылка $\mathrm{Prog}(X)$ при данном $x$, так что немедленно $x\in X$.

\emph{$\mathrm{Ind}(Y)$.} \textit{База:} по (O1) $[0,0)=\varnothing\subseteq X$, так что $0\in Y$. \textit{Шаг:} пусть $z\in Y$, то есть $[0,z)\subseteq X$; по уже доказанному $Y\subseteq X$ также $z\in X$, так что $[0,z)\cup\{z\}\subseteq X$. По (O2), $[0,Sz)=[0,z)\cup\{z\}$ (ведь $y<Sz \leftrightarrow y\leqslant z \leftrightarrow y<z\,\lor\,y=z$), так что $[0,Sz)\subseteq X$, то есть $Sz\in Y$.

Применяя $I$ к $Y$, получаем $\mathrm{Tot}(Y)$. Вместе с уже установленным $Y\subseteq X$ это немедленно даёт $\mathrm{Tot}(X)$ -- заключение $I^*$.
\end{proof}

Асимметрия направлений здесь содержательна. Направление $I^*\Rightarrow I$ (Лемма выше) существенно использует все три факта (O1)--(O3), причём именно (O3) играет решающую роль -- в разборе случаев <<$x$ есть $0$ или последователь>>. Направление $I\Rightarrow I^*$, напротив, обходится фактами (O1)--(O2) и в его доказательстве (O3) попросту не используется.

\textbf{(O3) существенно: модель, где $I^*$ есть, а $I$ -- нет.} Возьмём в качестве носителя ординал $\omega+\omega$ -- две последовательные копии $\mathbb N$, с обычным ординальным порядком $<$, и определим последователь $S$ \emph{отдельно внутри каждой копии}: $S(n)=n+1$ в первой копии, $S(\omega+n)=\omega+n+1$ во второй. Функция $S$ всюду определена и инъективна, (O1)--(O2) выполнены, но элемент $\omega$ (начало второй копии) не есть $S(y)$ ни для какого $y$ -- это ровно предельный ординал, на котором ломается (O3).

Поскольку $\omega+\omega$ -- честный ординал, то есть вполне упорядоченное множество, трансфинитная индукция (§1.3) верна в нём для \emph{произвольных} подмножеств носителя -- а не только для подмножеств, задаваемых формулой, -- так что $I^*$ (а вместе с ним $L$ и $D$) выполняется здесь в подлинно втором порядке, без всяких оговорок про схему.\footnote{Замечание для осторожного читателя: не значит ли это, что индукция вообще <<навязывает>> модели быть вполне упорядоченной, то есть стандартной? Нет: существуют нестандартные модели $\PA$, в которых полная схема $I$ (а с ней и $I^*,L,D$ как схемы) выполнена, однако сама модель, если смотреть на неё \emph{извне}, вполне упорядоченной не является -- в ней есть (неопределимые в языке $\PA$) подмножества без наименьшего элемента. Противоречия здесь нет: как схема аксиом, $I^*$ гарантирует наименьший элемент лишь для подмножеств, \emph{определимых} формулой языка $\PA$; на «неправильные» подмножества нестандартной модели эта схема просто не распространяется. В примере выше мы законно потребовали $I^*$ для абсолютно всех подмножеств носителя именно потому, что $\omega+\omega$ по построению стандартна (совпадает с настоящим ординалом); для нестандартных моделей такой роскоши уже нет.} А вот $I$ -- нет: множество $X$ = первая копия $\mathbb N$ содержит $0$ и замкнуто относительно $S$ (то есть $\mathrm{Ind}(X)$), но $X\ne\omega+\omega$. Именно это и предсказывает асимметрия выше: модель нарушает (O3) -- и вместе с ним $I$, -- но никак не трогает чисто логическую тройку $I^*, L, D$.

\subsection{Схема коллекции $B$ и индукция}
Природа связи между схемой коллекции $B$ и схемой индукции $I$ иная и заметно более тонкая метаматематически, чем всё рассмотренное выше: она не сводится ни к чистой логике, ни к минимальной арифметике последователя. Мы не будем доказывать её здесь, ограничившись точной формулировкой известного результата.

Внутри слабой базовой теории $\PA^-$ (аксиомы дискретного упорядоченного полукольца без индукции, см. Kaye, \textit{Models of Peano Arithmetic}) схема $B$ сама по себе \textbf{не влечёт} схему $I$. Но уже относительно базовой индукции по $\Delta_0$-формулам ($I\Delta_0$) связь восстанавливается и принимает вид строгой иерархии (Paris--Kirby, 1978):
$$I\Sigma_n \;\toinf\; B\Sigma_n \;\toinf\; I\Sigma_{n-1} \qquad (n\ge1),$$
а на уровне полных (неограниченных по сложности формулы) схем -- $I \iff B$ относительно базы $I\Delta_0$. Иерархия связей между всеми пятью принципами показана на диаграмме:

\begin{center}
\catcode`\"=12
\begin{tikzcd}[row sep=3em, column sep=4.5em]
|[draw, rounded corners, inner sep=8pt]| I^* \Leftrightarrow L \Leftrightarrow D
\arrow[r, shift left=2, Rightarrow, "\text{(O2),(O3)}"]
& |[draw, rounded corners, inner sep=6pt]| I
\arrow[l, shift left=2, Rightarrow, "\text{(O1),(O2)}"']
\arrow[r, shift left=2, Rightarrow, "I\Delta_0"]
& |[draw, rectangle, rounded corners, inner sep=6pt]| B
\arrow[l, shift left=2, Rightarrow, "I\Delta_0"']
\end{tikzcd}
\end{center}

Равносильность $I^*\Leftrightarrow L\Leftrightarrow D$ доказана в §2.2 чисто логически, с точностью до трихотомии $<$. Переход к $I$ требует минимальной арифметики последователя (§2.3): именно аксиома (O3) создаёт асимметрию -- она нужна для $I^*\Rightarrow I$, но не для $I\Rightarrow I^*$. Переход к схеме коллекции $B$ существенно использует индукцию $I\Delta_0$; доказательство не приводится.

\section{Следствия принципа индукции в математических структурах}

Ниже приведен перечень глубоких математических утверждений, справедливость которых существенно завязана на дедуктивную силу индукции. Там, где это возможно, указан точный статус утверждения в терминах подсистем арифметики второго порядка ($\textsf{RCA}_0$, $\textsf{WKL}_0$, $\textsf{ACA}_0$).

\subsection{Принцип Дирихле (Pigeonhole Principle)}
Функциональная формулировка: если $A$ и $B$ --- конечные множества, причем мощность $A$ строго больше мощности $B$, то не существует инъективного отображения $f: A \toinf B$.
\textit{Статус:} Доказуемость принципа Дирихле является классическим вопросом теории доказательств. По теореме Пэриса--Уилки--Вудса и Айтаи (1988), $PHP$ даже для $\Delta_0$-определимых функций не доказуем в $I\Delta_0$ --- то есть одна из самых элементарных комбинаторных истин невыводима из ограниченной индукции, что демонстрирует крайнюю слабость $I\Delta_0$ как базовой теории.

\subsection{Лемма Кёнига о бесконечном пути}
Формулировка в языке частично упорядоченных множеств: Если конечно-ветвящееся частично упорядоченное множество $(T, \leqslant_T)$, являющееся деревом с единственным корнем, содержит бесконечное число элементов, то оно содержит бесконечную цепь (путь).
\textit{Статус:} Слабая лемма Кёнига ($\textsf{WKL}_0$, для бинарных деревьев) эквивалентна компактности отрезка в матанализе. Полная лемма Кёнига (для любых конечно-ветвящихся деревьев) строго мощнее и эквивалентна арифметической схеме свертывания $\textsf{ACA}_0$. Дополнительная сила по сравнению с $\textsf{WKL}_0$ возникает не из индукции (схема индукции в обеих системах одна и та же), а именно из арифметической схемы свертывания: чтобы на каждом шаге выбрать поддерево с бесконечным числом вершин, нужно разрешить предикат «поддерево бесконечно», который для произвольных (не только бинарных) деревьев не сводится к $\Delta_0$-условию и требует именно $\textsf{ACA}_0$.

\subsection{Алгебраические законы в $\PA$}
Коммутативность и ассоциативность сложения ($\forall x, y \; x+y = y+x$) и умножения.

\textit{Уточнение базовой теории.} Здесь важно оговорить, о какой именно безындукционной теории идёт речь. Теория $\PA^-$, использованная в §2.4 (теория неотрицательных частей дискретно упорядоченного кольца, см. Kaye, \textit{Models of Peano Arithmetic}), уже \emph{постулирует} коммутативность и ассоциативность $+$ и $\times$ как отдельные неиндукционные аксиомы --- поэтому о «некоммутативных моделях $\PA^-$» говорить нельзя: это противоречило бы самому определению $\PA^-$. Утверждение ниже относится к более слабой теории $\PArec$, введённой в §2.3, -- арифметике Робинсона в её исходной, рекурсивной форме, с явными аксиомами порядка (O1)--(O3), но без индукции и без отдельного постулирования коммутативности или ассоциативности.

\textit{Модель без коммутативности:} Вывести переместительный закон в $\PArec$ без индукции невозможно: существуют (необходимо нестандартные) модели $\PArec$, в которых сложение некоммутативно. Содержательная иллюстрация того же явления --- хотя и не буквальная модель $\PArec$, а структура другого типа (упорядоченный моноид) --- арифметика ординалов: для первого трансфинитного ординала $\omega$ справедливо $\omega + 1 \neq 1 + \omega$. Именно невозможность вывести коммутативность без индукции и есть причина, по которой $\PA^-$ в смысле Kaye включает коммутативность и ассоциативность в список аксиом явно, а не пытается получить их как следствия.

\subsection{Теория конечных множеств по Дедекинду}
\begin{enumerate}[label=\arabic*.]
    \item \textbf{Собственные подмножества:} Утверждение «Если $B \subsetneq A$ и $A$ конечно, то мощность $B$ строго меньше мощности $A$» является прямым следствием индукции по числу элементов множества $A$. Без индукции конечные множества ведут себя как квазибесконечные.
    \item \textbf{Конечность по Дедекинду:} Множество $A$ конечно по Дедекинду, если не существует инъекции $f: A \toinf A$, не являющейся сюръекцией. С помощью индукции легко доказывается направление «стандартно конечно $\Rightarrow$ конечно по Дедекинду». Обратное направление, однако, из одной индукции не следует: оно доказуемо с помощью счётной аксиомы выбора $AC_\omega$ (если множество бесконечно, $AC_\omega$ позволяет выбрать в нём счётную подпоследовательность различных элементов, что и даёт нужную несюръективную инъекцию) -- но это лишь \emph{следствие} $AC_\omega$, а не равносильное ей утверждение. В любом случае, без всякого выбора в теории $\textsf{ZF}$ это направление недоказуемо: существуют модели $\textsf{ZF}$ с бесконечными дедекиндово-конечными множествами (модели Френкеля--Мостовского, первая модель Коэна).
\end{enumerate}

\subsection{Теорема Рамсея ($RT^2_2$)}
Теорему Рамсея для пары цветов понятнее всего формулировать на языке обычных графов, без разговора о раскрасках: для двух цветов один из них -- это <<есть ребро>>, а другой -- <<ребра нет>>, так что произвольная раскраска пар вершин в 2 цвета -- это то же самое, что произвольный граф на этих вершинах.

\textbf{Конечный случай.} Классический пример -- <<теорема о вечеринке>>: среди любых 6 человек (рассмотрим граф знакомств, где ребро -- это знакомство) найдутся либо трое, попарно знакомых друг с другом (треугольник в графе), либо трое, попарно не знакомых между собой (треугольник в его дополнении); 6 -- наименьшее такое число (число Рамсея $R(3,3)=6$). В общем виде: для любых $k,m$ существует число Рамсея $R(k,m)$ -- наименьшее $N$, при котором любой граф на $N$ вершинах содержит либо клику размера $k$, либо независимое множество размера $m$. Стоит заметить, что даже это чисто конечное по своему содержанию утверждение (лишь квантифицированное по $k,m$) доказывается индукцией, по сумме $k+m$.

\textbf{Бесконечный случай ($RT^2_2$).} В любом графе на бесконечном множестве вершин найдётся бесконечное подмножество вершин, образующее либо полный подграф (все рёбра есть), либо вполне несвязный подграф (рёбер нет вовсе). Это прямое следствие индукции (точнее, принципа $L$: среди <<типов>> вершин относительно уже отобранных всегда найдётся бесконечно много одинаковых).\footnote{Точная обратно-математическая сила $RT^2_2$ известна и весьма тонка: она доказуема в $\textsf{RCA}_0+I\Sigma_2$, но сама не влечёт $I\Sigma_2$ (Чолак--Джокуш--Слейман, 2001) -- мы не будем углубляться в эти детали.}

\subsection{Тотальность функции Аккермана}
По теореме Парсонса, доказуемо рекурсивные функции $I\Sigma_1$ суть в точности примитивно рекурсивные функции; поскольку функция Аккермана по построению не примитивно рекурсивна, ее тотальность \textbf{недоказуема} в $I\Sigma_1$. Она становится доказуемой уже на следующем уровне, в $I\Sigma_2$ (что соответствует месту функции Аккермана в быстрорастущей иерархии, на уровне $\omega$, тогда как $I\Sigma_1$ покрывает лишь конечные уровни этой иерархии).

\subsection{Пределы индукции: теорема Гудстейна и принцип Пэриса--Харрингтона}
Все примеры §§3.1--3.6 показывали, сколь многого можно добиться средствами индукции $I$. Два следующих результата демонстрируют прямо противоположное: у индукции $I$ (то есть у полной схемы индукции $\PA$) есть чёткий, точно измеримый \textbf{предел} мощности -- и именно эти два примера его маркируют.

\textbf{Теорема Гудстейна.} Сформулирована и доказана (безо всякой связи с метаматематикой $\PA$) Р.\,Гудстейном ещё в 1944 году: для любого числа $n$ последовательность Гудстейна $g_n(0), g_n(1), g_n(2),\ldots$ (определяемая через запись чисел в наследственной позиционной системе по основанию $k$, с заменой основания $k$ на $k{+}1$ и вычитанием единицы на каждом шаге) всегда достигает $0$ за конечное число шагов. Это истинное $\Pi_2$-предложение языка $\PA$ (вида $\forall n\,\exists k\,g_n(k)=0$); лишь спустя почти четыре десятилетия Кирби и Пэрис (1982) показали, что оно \textbf{не выводимо в $\PA$}.

\textbf{Принцип Пэриса--Харрингтона.} Усиленный вариант конечной теоремы Рамсея из §3.5 (с дополнительным требованием <<относительной большости>> искомого монохроматического множества) -- истинное, чисто комбинаторное предложение (Пэрис, Харрингтон, 1977), которое было \emph{сразу же}, при самом своём появлении, предъявлено как пример недоказуемого в $\PA$ утверждения. В этом смысле принцип Пэриса--Харрингтона -- первый пример \emph{натуральной} (не метаматематической по формулировке, не апеллирующей явно к доказуемости или к гёделевской диагонализации) теоремы, для которой недоказуемость в $\PA$ была установлена одновременно с самой теоремой: теорема Гудстейна как \emph{математическое утверждение} появилась значительно раньше (1944), но её независимость от $\PA$ была обнаружена лишь позже -- под влиянием результата Пэриса и Харрингтона.

Оба факта тем не менее доказуемы средствами, выходящими за пределы $I$, -- в частности, уже в $\ZFC$ и в полной второпорядковой арифметике $\Ztwo$: достаточно \emph{трансфинитной} индукции вдоль ординала $\varepsilon_0$, доступной уже в относительно скромной подсистеме $\textsf{ATR}_0\subseteq\Ztwo$. Ординал $\varepsilon_0$ появляется здесь не случайно: по теореме Генцена (1936), это в точности доказательно-теоретический ординал $\PA$ -- граница, которую индукция $I$ не может формализовать как схему доказательства сама по себе. Иными словами: локальная и возвратная схемы индукции $I\iff I^*$ -- при всей демонстрируемой ими мощи -- сами по себе \textbf{не замкнуты} относительно всех истинных арифметических утверждений; выход за их пределы требует принципиально более сильной формы индукции -- трансфинитной, -- недоступной как первопорядковая схема аксиом над стандартной моделью $\omega$.

\section*{Список обозначений}
\begin{description}[leftmargin=3.2cm, style=nextline, font=\normalfont]
    \item[$\PA$] Арифметика Пеано (первопорядковая, с полной схемой индукции $I$).
    \item[$\Ztwo$] Арифметика второго порядка (анализ) с непредикативной схемой свёртывания.
    \item[$\ZFC$] Теория множеств Цермело--Френкеля с аксиомой выбора.
    \item[$\PA^-$] Аксиомы дискретного упорядоченного полукольца без индукции (в смысле Kaye), с явно постулированными коммутативностью и ассоциативностью $+,\times$.
    \item[$\PArec$] Арифметика Робинсона в рекурсивной форме (уравнения для $+,\times,S$) с явными аксиомами порядка (O1)--(O3), без индукции и без отдельного постулирования коммутативности/ассоциативности (§2.3).
    \item[$S$, $<$, $\leqslant$] Последователь; строгий и нестрогий порядок на числах.
    \item[$[0,x)$] Начальный интервал $\{y \mid y<x\}$.
    \item[$\mathrm{Ind}(X)$] Индуктивность $X$: $0\in X\,\land\,\forall x\,(x\in X\toinf Sx\in X)$.
    \item[$\mathrm{Prog}(X)$] Прогрессивность $X$: $\forall x\,([0,x)\subseteq X\toinf x\in X)$.
    \item[$\mathrm{Tot}(X)$] Тотальность $X$: $\forall x\,(x\in X)$.
    \item[$I$] Локальная индукция: $\forall X(\mathrm{Ind}(X)\toinf\mathrm{Tot}(X))$.
    \item[$I^*$] Возвратная (полная, course-of-values) индукция: $\forall X(\mathrm{Prog}(X)\toinf\mathrm{Tot}(X))$.
    \item[$L$] Принцип наименьшего элемента.
    \item[$D$] Принцип бесконечного спуска (отсутствия бесконечно убывающих цепей).
    \item[$B$] Схема коллекции (collection scheme).
    \item[(O1)--(O3)] Аксиомы $\PArec$ о согласовании $<$ и $S$: $0$ -- наименьший; $y<Sx\leftrightarrow y\leqslant x$; всякий ненулевой элемент есть последователь.
    \item[$I\Sigma_n$, $B\Sigma_n$, $I\Delta_0$] Индукция (соотв. коллекция) для формул ограниченной кванторной сложности; иерархия подсистем $\PA$.
    \item[$\textsf{RCA}_0$] Базовая система обратной математики: рекурсивная схема свёртывания плюс $I\Sigma_1$.
    \item[$\textsf{WKL}_0$] $\textsf{RCA}_0$ плюс слабая лемма Кёнига (для бинарных деревьев).
    \item[$\textsf{ACA}_0$] $\textsf{RCA}_0$ плюс арифметическая схема свёртывания.
    \item[$\textsf{ATR}_0$] $\textsf{RCA}_0$ плюс арифметическая трансфинитная рекурсия; строго сильнее $\textsf{ACA}_0$.
    \item[$\varepsilon_0$] Доказательно-теоретический ординал $\PA$ (Генцен, 1936): супремум порядковых типов канонических (натурально заданных) вполне упорядоченных отношений на $\mathbb N$, фундированность которых доказуема в $\PA$.\footnote{Не наименьший ординал, для которого $\PA$ не может доказать фундированность: существуют куда более короткие, но искусственно (через предикат непротиворечивости) устроенные порядки с тем же дефектом. Оговорка про <<канонические>> порядки как раз и нужна, чтобы отсечь такие патологии -- вопрос о точной границе этого понятия является предметом отдельной дискуссии в теории доказательств.}
    \item[$RT^2_2$] Теорема Рамсея для пар и двух цветов (в графовой форме -- для произвольного графа на бесконечном множестве вершин).
\end{description}

\end{document}